\documentclass[12pt, a4paper]{article}

% ─── Packages ────────────────────────────────────────────────────────────────
\usepackage[utf8]{inputenc}
\usepackage[T1]{fontenc}
\usepackage{amsmath, amssymb, amsthm}
\usepackage{mathtools}
\usepackage{geometry}
\usepackage{hyperref}
\usepackage{booktabs}
\usepackage{array}
\usepackage{enumitem}
\usepackage{xcolor}
\usepackage{titlesec}
\usepackage{parskip}
\usepackage{fancyhdr}
\usepackage{graphicx}
\usepackage{tikz}
\usetikzlibrary{arrows.meta, positioning, shapes.geometric}

\geometry{margin=2.5cm, headheight=15pt}
\hypersetup{colorlinks=true, linkcolor=blue!60!black, citecolor=blue!60!black, urlcolor=blue!60!black}

% ─── Theorem environments ─────────────────────────────────────────────────────
\theoremstyle{plain}
\newtheorem{theorem}{Theorem}[section]
\newtheorem{lemma}[theorem]{Lemma}
\newtheorem{proposition}[theorem]{Proposition}
\newtheorem{corollary}[theorem]{Corollary}

\theoremstyle{definition}
\newtheorem{definition}[theorem]{Definition}
\newtheorem{example}[theorem]{Example}
\newtheorem{axiom}[theorem]{Axiom}

\theoremstyle{remark}
\newtheorem{remark}[theorem]{Remark}
\newtheorem{claim}[theorem]{Claim}

% ─── Notation macros ──────────────────────────────────────────────────────────
\newcommand{\ket}[1]{\left| #1 \right\rangle}
\newcommand{\bra}[1]{\left\langle #1 \right|}
\newcommand{\braket}[2]{\left\langle #1 \mid #2 \right\rangle}
\newcommand{\norm}[1]{\left\| #1 \right\|}
% \Psi is already defined by LaTeX
\newcommand{\pspace}{\ensuremath{\mathsf{PSPACE}}}
\newcommand{\bqp}{\ensuremath{\mathsf{BQP}}}
\newcommand{\bpp}{\ensuremath{\mathsf{BPP}}}
\newcommand{\pp}{\ensuremath{\mathsf{PP}}}
\newcommand{\Pcls}{\ensuremath{\mathsf{P}}}
\newcommand{\NP}{\ensuremath{\mathsf{NP}}}
\newcommand{\DAG}{\ensuremath{\mathrm{DAG}}}
\newcommand{\RWX}{\ensuremath{\mathrm{RWX}}}
\newcommand{\SD}{\ensuremath{\mathrm{SD}}}   % Superposition Decay
\newcommand{\QC}{\ensuremath{\mathcal{Q}}}   % Quantum circuit family
\newcommand{\Nat}{\mathbb{N}}
\newcommand{\Real}{\mathbb{R}}
\newcommand{\Complex}{\mathbb{C}}

% ─── Header / Footer ─────────────────────────────────────────────────────────
\pagestyle{fancy}
\fancyhf{}
\rhead{\small OBINexus Theoretical Computing}
\lhead{\small Quantum Superposition Decay}
\cfoot{\thepage}

% ─────────────────────────────────────────────────────────────────────────────
\begin{document}

% ─── Title ───────────────────────────────────────────────────────────────────
\begin{titlepage}
  \centering
  \vspace*{2cm}
  {\LARGE \textbf{Quantum Superposition Decay for Quantum Complexity Classes}}\\[1em]
  {\large A Classical--Quantum Bridge via the Read-Write-Execute Memory Model}\\[2em]
  {\large Nnamdi Okpala}\\[0.5em]
  {\normalsize OBINexus Computing}\\[0.5em]
  {\normalsize \texttt{okpalan@protonmail.com}}\\[2em]
  {\normalsize \today}\\[3em]

  \begin{abstract}
    We introduce \emph{Superposition Decay} (\SD) as a formal measurement
    window that governs the collapse of a quantum superposition state on a
    classical computing substrate.  Building on standard quantum postulates
    and the wave-function collapse principle, we show that any classical
    algorithm bounded by linear time complexity---deterministic or
    non-deterministic---can be evaluated in parallel via superposition decay
    on a quantum system.  We further demonstrate that all problems belonging
    to a system of linear equations can be expressed as a probabilistic
    directed acyclic graph (\DAG) on such a system.  A concrete
    Read-Write-Execute (\RWX) memory model formalises the decay window in
    operational terms: four write operations and two read operations reduce
    atomically to a single coherent read.  We map these results onto the
    standard quantum complexity hierarchy ($\Pcls \subseteq \bpp \subseteq
    \bqp \subseteq \pp \subseteq \pspace$) and characterise which
    promise-problem classes are tractable under superposition decay.  The
    framework is then extended in two directions: (i) \textbf{Rift}, a
    quantum-classical grammar matcher whose parser exploits grammar
    superposition to evaluate all production rules in parallel, collapsing
    to a unique parse tree at a grammar decay event; and (ii) the
    \textbf{Quantum Burst Protocol (QBP)}, a communication protocol in which
    all candidate network paths are held in superposition during handshake
    and selected by a single measurement event.  Both components are
    implemented within the MMUKO OS / OBINexus toolchain via the pipeline
    \texttt{riftlang.exe} $\to$ \texttt{.so.a} $\to$ \texttt{rift.exe} $\to$
    \texttt{gosilang}, orchestrated by \texttt{nlink} and \texttt{polybuild}.
  \end{abstract}
  \vfill
\end{titlepage}

\tableofcontents
\newpage

% ─────────────────────────────────────────────────────────────────────────────
\section{Introduction}
\label{sec:intro}
% ─────────────────────────────────────────────────────────────────────────────

The relationship between quantum and classical computation is most precisely
characterised through the lens of promise-problem complexity classes.
While quantum speedup for specific problem families (integer factorisation,
unstructured search) is well established, a general framework that maps the
\emph{measurement event}---the moment a quantum superposition collapses to a
definite classical outcome---onto classical computational primitives has
received less systematic attention.

This paper proposes such a framework under the name \textbf{Superposition
Decay} (\SD).  The central intuition is that the collapse of a superposition
is not a passive observation but an active computational window: a bounded
interval of control during which a system may coherently commit or roll back
parallel computations, analogous to a transactional memory operation.

We formalise five related contributions:

\begin{enumerate}[label=\textbf{C\arabic*.}]
  \item \textbf{The \SD{} window.}  A formal definition of the measurement
        window as a triple (resolved / rejected / pending), bounding
        acceptance and rejection by probability thresholds $a(n)$ and $b(n)$,
        with an intermediate \emph{pending} region supporting deferred
        evaluation.

  \item \textbf{The RWX memory model.}  A Read-Write-Execute abstraction that
        maps quantum measurement to classical memory transactions: $n$ writes
        followed by $k$ reads collapse to a single coherent read event when
        the last read is the terminal operation.

  \item \textbf{Complexity-class tractability under \SD.}  A characterisation
        showing that $\bqp$ promise problems are exactly the class tractable
        under a single superposition decay event; problems in $\pp$ and
        $\pspace$ require multiple decay events or additional measurement
        principles.

  \item \textbf{Rift: quantum-classical grammar matching.}  A formalisation
        of the Rift compiler's parser as a quantum automaton, with grammar
        production rules as superposed branches and parse commitment as a
        grammar decay event governed by the stabilisation principle.

  \item \textbf{Quantum Burst Protocol (QBP).}  A communication protocol in
        which all candidate network paths are held in superposition during
        connection negotiation and selected by a single measurement (spectral
        decay), achieving optimal routing in $O(1)$ measurement time.
\end{enumerate}

The paper is organised as follows.  Section~\ref{sec:prelim} reviews necessary
background in quantum mechanics and complexity theory.  Section~\ref{sec:sd}
defines superposition decay formally.  Section~\ref{sec:rwx} presents the
RWX memory model.  Section~\ref{sec:complexity} maps results onto the
complexity hierarchy.  Section~\ref{sec:parallel} states and proves the main
parallelism theorem.  Section~\ref{sec:rift} formalises Rift as a
quantum-classical grammar matcher.  Section~\ref{sec:qbp} introduces QBP.
Section~\ref{sec:mmuko} describes the MMUKO OS and OBINexus toolchain.
Section~\ref{sec:conclusion} concludes.

% ─────────────────────────────────────────────────────────────────────────────
\section{Preliminaries}
\label{sec:prelim}
% ─────────────────────────────────────────────────────────────────────────────

\subsection{Quantum Postulates}

We work in the standard Hilbert-space formulation of quantum mechanics.  A
quantum state is represented by a unit vector $\ket{\psi} \in \Complex^{2^n}$
for an $n$-qubit system.

\begin{axiom}[State Postulate]
  The complete state of a quantum system at time $t$ is described by a
  normalised vector $\ket{\psi(t)} \in \mathcal{H}$, the system's Hilbert
  space, satisfying $\braket{\psi(t)}{\psi(t)} = 1$.
\end{axiom}

\begin{axiom}[Superposition Postulate]
  If $\ket{\psi_1}$ and $\ket{\psi_2}$ are valid states of a quantum system,
  then any linear combination
  \[
    \ket{\psi} = c_1 \ket{\psi_1} + c_2 \ket{\psi_2},
    \quad c_1, c_2 \in \Complex, \quad |c_1|^2 + |c_2|^2 = 1,
  \]
  is also a valid state.  This follows from the linearity and homogeneity of
  the Schr\"{o}dinger equation.
\end{axiom}

\begin{axiom}[Wave-Function Collapse / Measurement Postulate]
  Upon measurement of an observable $\hat{O}$ with eigenvalues
  $\{\lambda_i\}$ and eigenstates $\{\ket{e_i}\}$, the system transitions
  from $\ket{\psi} = \sum_i c_i \ket{e_i}$ to eigenstate $\ket{e_j}$ with
  probability $|c_j|^2$.  The post-measurement state is $\ket{e_j}$.
\end{axiom}

\begin{axiom}[Evolution Postulate]
  Between measurements, a closed quantum system evolves unitarily:
  $\ket{\psi(t)} = U(t)\ket{\psi(0)}$, where $U(t)$ is a unitary operator.
\end{axiom}

\subsection{Classical Turing Machine and Complexity}

A \emph{Turing machine} (TM) is a tuple $M = (Q, \Sigma, \Gamma, \delta,
q_0, q_{\text{acc}}, q_{\text{rej}})$ with the standard interpretation.
Time complexity $T(n)$ denotes the maximum number of steps on inputs of
length $n$.  We use the standard complexity classes $\Pcls$, $\NP$, $\bpp$,
$\pp$, $\pspace$, and $\bqp$; their promise-problem definitions are recalled
in Section~\ref{sec:complexity}.

\subsection{Quantum Circuits}

A \emph{quantum circuit} $Q_n$ on $n$ qubits is a sequence of gates drawn
from a universal gate set (e.g., $\{H, \mathrm{CNOT}, T\}$) that accepts an
$n$-qubit input and produces a 1-qubit output.  A \emph{uniform circuit
family} $\QC = \{Q_n : n \in \Nat\}$ is one where the description of $Q_n$
is computable in polynomial time from $n$.

\subsection{Probabilistic Directed Acyclic Graphs}

\begin{definition}[Probabilistic DAG]
  A \emph{probabilistic directed acyclic graph} $G = (V, E, p)$ is a DAG
  where each edge $(u, v) \in E$ carries a probability weight
  $p(u,v) \in [0,1]$ such that for every vertex $u$,
  $\sum_{v : (u,v) \in E} p(u,v) \leq 1$.  Each path from a source to a
  sink represents a computation branch; the probability of a path is the
  product of its edge weights.
\end{definition}

% ─────────────────────────────────────────────────────────────────────────────
\section{Superposition Decay}
\label{sec:sd}
% ─────────────────────────────────────────────────────────────────────────────

\subsection{Informal Motivation}

When a quantum system in superposition is measured, the resulting classical
outcome is not predetermined---it is drawn from a probability distribution
defined by the amplitudes.  We argue that this \emph{measurement event} is
best understood not as an instantaneous point but as a \emph{window}: a
bounded region of the computation during which the system retains coherence
and the outcomes of parallel branches may still be combined or discarded.

Once the window closes---the superposition \emph{decays}---a single classical
outcome is committed.  The \SD{} framework captures this window operationally.

\subsection{Formal Definition}

\begin{definition}[Quantum Connection Graph]
  A \emph{quantum connection} is a graph $G = (E, N)$ where $N$ is the set
  of $n$-qubit nodes (quantum registers) and $E \subseteq N \times N$ is the
  set of entanglement edges.  A path from node $A$ to node $B$ in $G$
  represents a sequence of entangled operations connecting two computational
  loci.
\end{definition}

\begin{definition}[Superposition State Vector]
  For an $n$-qubit system, the superposition state is
  \[
    \ket{\Psi} = \sum_{x \in \{0,1\}^n} \alpha_x \ket{x},
    \quad \alpha_x \in \Complex, \quad \sum_{x} |\alpha_x|^2 = 1.
  \]
  The probability of observing basis state $\ket{x}$ is $p(x) = |\alpha_x|^2$.
\end{definition}

\begin{definition}[Superposition Decay Window]
  \label{def:sdwindow}
  Let $\QC = \{Q_n\}$ be a uniform quantum circuit family and let
  $a, b : \Nat \to [0,1]$ be functions with $a(n) > b(n)$ for all $n$.
  The \emph{superposition decay window} $\SD(a, b)$ is the triple
  \[
    \SD(a, b) \;=\; \bigl(\, \text{resolved},\; \text{rejected},\; \text{pending} \,\bigr),
  \]
  where for an input $x$ of length $n$:
  \begin{itemize}
    \item $x$ is \textbf{resolved} (accepted) if $\Pr[Q_{|x|}(x) = 1] \geq a(|x|)$,
    \item $x$ is \textbf{rejected} if $\Pr[Q_{|x|}(x) = 1] \leq b(|x|)$,
    \item $x$ is \textbf{pending} if $b(|x|) < \Pr[Q_{|x|}(x) = 1] < a(|x|)$.
  \end{itemize}
  The \emph{decay event} is the transition of a state from \emph{pending} to
  either \emph{resolved} or \emph{rejected} upon measurement.
\end{definition}

\begin{remark}
  The three-valued output (resolved / rejected / pending) generalises the
  standard binary outcome of a classical TM.  The \emph{pending} state
  corresponds precisely to the undecided region of a promise problem: the
  system is contractually guaranteed not to receive inputs from this region
  in the promise-problem formulation, but the \SD{} framework makes the
  region explicit as a computational state.
\end{remark}

\subsection{Constructive and Destructive Superposition}

The superposition principle admits both \emph{constructive} interference
(amplitudes add, increasing probability of a desired outcome) and
\emph{destructive} interference (amplitudes cancel, suppressing undesired
outcomes).  We formalise this in terms of the wave function as follows.

\begin{definition}[Constructive / Destructive Interference]
  Given two states $\ket{\psi_1} = \sum_x \alpha_x \ket{x}$ and
  $\ket{\psi_2} = \sum_x \beta_x \ket{x}$, their superposition
  $\ket{\psi} = c_1 \ket{\psi_1} + c_2 \ket{\psi_2}$ exhibits:
  \begin{itemize}
    \item \textbf{Constructive interference} at $x$ if
          $|c_1 \alpha_x + c_2 \beta_x| > \max(|c_1 \alpha_x|, |c_2 \beta_x|)$.
    \item \textbf{Destructive interference} at $x$ if
          $|c_1 \alpha_x + c_2 \beta_x| < \min(|c_1 \alpha_x|, |c_2 \beta_x|)$.
  \end{itemize}
\end{definition}

In the \SD{} framework, a quantum algorithm exploits constructive interference
to amplify the amplitude of the correct solution branch before the decay
window closes, and destructive interference to suppress incorrect branches.

% ─────────────────────────────────────────────────────────────────────────────
\section{The Read-Write-Execute (RWX) Memory Model}
\label{sec:rwx}
% ─────────────────────────────────────────────────────────────────────────────

\subsection{Motivation}

Classical memory access is characterised by read ($R$), write ($W$), and
execute ($X$) permissions, typically represented as a triple $(r, w, x)$ per
memory region.  We propose a quantum analogue in which the \emph{sequence} of
reads and writes within a superposition decay window determines the atomicity
and coherence of the resulting classical observation.

\subsection{Formal Model}

\begin{definition}[RWX Operation Sequence]
  An \emph{RWX sequence} is a finite ordered list of operations
  $\sigma = (o_1, o_2, \ldots, o_k)$ where each $o_i \in \{R, W, X\}$.
  The sequence is \emph{terminal-read} if $o_k = R$ (the final operation is
  a read).
\end{definition}

\begin{definition}[Atomic Decay Unit]
  \label{def:adu}
  A terminal-read RWX sequence of $w$ writes followed by $r$ reads,
  \[
    \sigma = (\underbrace{W, W, \ldots, W}_{w},\;
              \underbrace{R, R, \ldots, R}_{r-1},\; R),
  \]
  constitutes a single \emph{atomic decay unit} (ADU) equivalent to one
  coherent read operation.  Formally:
  \[
    w \cdot W + r \cdot R \;\xrightarrow{\;\SD\;}\; 1 \cdot R_{\text{coherent}}.
  \]
\end{definition}

\begin{example}[4-Write 2-Read ADU]
  The concrete instance stated in the source document is:
  \[
    4W + 2R \;\xrightarrow{\;\SD\;}\; 1R.
  \]
  This means: four parallel write operations (corresponding to four parallel
  computation branches in superposition) and two intermediate read probes
  collapse---upon superposition decay---into a single coherent classical read
  outcome.  The dual-state nature of the system (write = superposed,
  read = collapsed) ensures that the terminal read is the only externally
  observable event.
\end{example}

\begin{proposition}[Dual-State Equivalence]
  In a quantum system respecting the \SD{} window, the state before the
  terminal read is a superposition of all write branches.  The terminal read
  collapses this superposition to a single classical value.  Therefore, the
  external observer sees exactly one read operation regardless of the number
  of internal write branches.
\end{proposition}

\begin{proof}
  By the Measurement Postulate (Axiom~3), measurement of the system in state
  $\ket{\Psi} = \sum_x \alpha_x \ket{x}$ produces outcome $x$ with
  probability $|\alpha_x|^2$ and collapses the state to $\ket{x}$.  Each
  write branch corresponds to a basis state $\ket{x}$.  Since the terminal
  read is the sole measurement, all prior writes are unobserved intermediate
  computations within the unitary evolution $U(t)$.  By the Evolution
  Postulate (Axiom~4), $U(t)$ is reversible and does not produce classical
  outputs; only the terminal measurement does.  Hence the external observable
  record contains exactly one read event. \qed
\end{proof}

\subsection{RWX Permission Table}

The classical RWX permission system maps directly onto the decay window:

\begin{center}
\begin{tabular}{@{}lll@{}}
  \toprule
  \textbf{Permission} & \textbf{Classical Meaning} & \textbf{Quantum / \SD{} Analogue} \\
  \midrule
  $R$ (Read)    & Load value from memory         & Terminal measurement; collapses superposition \\
  $W$ (Write)   & Store value to memory          & Unitary gate application; maintains superposition \\
  $X$ (Execute) & Run code in memory region      & Quantum circuit evaluation within decay window \\
  \bottomrule
\end{tabular}
\end{center}

% ─────────────────────────────────────────────────────────────────────────────
\section{Complexity Classes Under Superposition Decay}
\label{sec:complexity}
% ─────────────────────────────────────────────────────────────────────────────

\subsection{Promise Problem Formulation}

A \emph{promise problem} is a pair $(A_{\text{yes}}, A_{\text{no}})$ of
disjoint languages.  An algorithm \emph{solves} the promise problem if it
accepts every $x \in A_{\text{yes}}$ and rejects every $x \in A_{\text{no}}$,
with no constraint on inputs outside $A_{\text{yes}} \cup A_{\text{no}}$.

\subsection{Complexity Class Definitions}

\begin{table}[h]
\centering
\renewcommand{\arraystretch}{1.5}
\begin{tabular}{@{}>{\bfseries}p{1.8cm}p{10.5cm}@{}}
  \toprule
  Class & Criteria \\
  \midrule
  $\Pcls$ &
    Promise problems for which a polynomial-time \emph{deterministic} TM
    accepts all strings in $A_{\text{yes}}$ and rejects all strings in
    $A_{\text{no}}$.
  \\
  $\bpp$ &
    Promise problems for which a polynomial-time \emph{probabilistic} TM
    accepts every string in $A_{\text{yes}}$ with probability $\geq \tfrac{2}{3}$
    and every string in $A_{\text{no}}$ with probability $\leq \tfrac{1}{3}$.
  \\
  $\bqp$ &
    Promise problems such that for functions $a, b : \Nat \to [0,1]$ with
    $a(n) \geq \tfrac{2}{3}$ and $b(n) \leq \tfrac{1}{3}$, there exists a
    uniform quantum circuit family $\QC = \{Q_n : n \in \Nat\}$ where
    $Q_n$ accepts $n$ qubits and gives a 1-qubit output, such that every
    $x \in A_{\text{yes}}$ is accepted with probability $\geq a(|x|)$ and
    every $x \in A_{\text{no}}$ is accepted with probability $\leq b(|x|)$.
  \\
  $\pp$ &
    Promise problems for which a polynomial-time probabilistic TM accepts
    every string in $A_{\text{yes}}$ with probability $> \tfrac{1}{2}$ and
    every string in $A_{\text{no}}$ with probability $\leq \tfrac{1}{2}$.
  \\
  $\pspace$ &
    Promise problems for which a \emph{deterministic} TM running in
    polynomial \emph{space} accepts all $A_{\text{yes}}$ and rejects all
    $A_{\text{no}}$.
  \\
  \bottomrule
\end{tabular}
\caption{Complexity class definitions in the promise-problem framework.}
\label{tab:classes}
\end{table}

\subsection{Tractability Under \SD}

\begin{theorem}[\SD-Tractability of \bqp]
  \label{thm:bqp}
  A promise problem $(A_{\text{yes}}, A_{\text{no}})$ is in $\bqp$ if and
  only if it is solvable by a uniform quantum circuit family within a single
  superposition decay window $\SD(a, b)$ with $a(n) - b(n) \geq \tfrac{1}{3}$
  for all $n$.
\end{theorem}

\begin{proof}
  $(\Rightarrow)$ By definition of $\bqp$, there exists a uniform circuit
  family $\{Q_n\}$ with the stated acceptance probabilities.  The circuit
  $Q_n$ operates entirely within a unitary evolution (all gates are write
  operations in RWX terms) and terminates with a single measurement (the
  terminal read).  This constitutes a single ADU as in
  Definition~\ref{def:adu}, hence falls within one \SD{} window.

  $(\Leftarrow)$ Suppose a single \SD$(a,b)$ window with $a(n) - b(n) \geq
  \tfrac{1}{3}$ solves the promise problem.  The circuit within the window is
  unitary (by the Evolution Postulate) and polynomial-time (by the
  polynomial-size constraint on the circuit family).  This is exactly the
  definition of a $\bqp$ circuit. \qed
\end{proof}

\begin{corollary}
  $\Pcls \subseteq \bpp \subseteq \bqp$ under the \SD{} framework, since
  deterministic and bounded-error probabilistic computations can be simulated
  by a single decay window with all amplitude concentrated on one branch.
\end{corollary}

\begin{theorem}[Multi-Decay Requirement for \pp]
  \label{thm:pp}
  There exist promise problems in $\pp \setminus \bqp$ that require more than
  one superposition decay event to resolve under \SD.
\end{theorem}

\begin{proof}[Proof Sketch]
  $\pp$ allows acceptance probabilities arbitrarily close to $\tfrac{1}{2}$,
  whereas $\bqp$ requires a gap of at least $\tfrac{1}{3}$.  A problem
  instance where $\Pr[\text{accept}] = \tfrac{1}{2} + \epsilon$ for
  arbitrarily small $\epsilon > 0$ cannot be distinguished from a reject
  instance (with $\Pr = \tfrac{1}{2} - \epsilon$) by a single decay event
  with fixed thresholds $a, b$ satisfying $a - b \geq \tfrac{1}{3}$.  Multiple
  sequential decay events (with amplitude amplification between events) are
  required to drive the probability above the \bqp{} threshold. \qed
\end{proof}

\begin{remark}
  $\pspace$ problems, being characterised by polynomial-space deterministic
  computation, lie outside the reach of a constant number of superposition
  decay events unless $\bqp = \pspace$, which is not currently believed.
  They require fundamentally different principles---either unbounded decay
  depth or oracle access---within the \SD{} framework.
\end{remark}

% ─────────────────────────────────────────────────────────────────────────────
\section{The Main Parallelism Theorem}
\label{sec:parallel}
% ─────────────────────────────────────────────────────────────────────────────

\subsection{Statement}

\begin{theorem}[Classical-to-Quantum Parallelism via \SD]
  \label{thm:main}
  Let $\mathcal{A}$ be a classical algorithm with time complexity $O(f(n))$
  where $f$ is a linear function ($f(n) = cn$ for some constant $c > 0$),
  and let $\mathcal{A}$ be deterministic or non-deterministic.  Then there
  exists a quantum circuit family $\QC_{\mathcal{A}} = \{Q_n\}$ such that:
  \begin{enumerate}
    \item $Q_n$ simulates all execution branches of $\mathcal{A}$ on inputs
          of length $n$ in \emph{parallel} within a single superposition decay
          window $\SD(a, b)$.
    \item The circuit size of $Q_n$ is polynomial in $n$.
    \item The correct output is obtained upon the terminal measurement (the
          decay event) with probability $\geq \tfrac{2}{3}$.
  \end{enumerate}
\end{theorem}

\begin{proof}
  We construct $Q_n$ as follows.

  \textbf{Step 1: Encode branches.}  The $k$ non-deterministic branches of
  $\mathcal{A}$ on input $x$ correspond to basis states $\ket{b_1}, \ldots,
  \ket{b_k}$ in a $\lceil \log_2 k \rceil$-qubit register.  Since
  $\mathcal{A}$ runs in linear time $O(n)$, the number of branches is at most
  $k \leq 2^{cn}$ for some constant $c$; encoding requires $O(n)$ qubits.

  \textbf{Step 2: Superpose branches.}  Apply a Hadamard transform $H^{\otimes
  \lceil \log_2 k \rceil}$ to create the uniform superposition
  $\ket{\Psi_0} = \frac{1}{\sqrt{k}} \sum_{i=1}^k \ket{b_i}$.

  \textbf{Step 3: Parallel evaluation.}  By the linearity of unitary
  evolution, the circuit $U_{\mathcal{A}}$ that implements one step of
  $\mathcal{A}$ acts on all branches simultaneously:
  $U_{\mathcal{A}} \ket{\Psi_0} = \frac{1}{\sqrt{k}} \sum_i U_{\mathcal{A}}
  \ket{b_i}$.  Since $\mathcal{A}$ runs in $O(n)$ steps, $U_{\mathcal{A}}$
  consists of $O(n)$ layers, each implementable by a polynomial-size circuit.

  \textbf{Step 4: Amplitude amplification.}  Apply Grover-type amplitude
  amplification~\cite{grover1996} to boost the amplitude of accepting branches.
  After $O(\sqrt{k})$ iterations, the probability of observing an accepting
  branch exceeds $\tfrac{2}{3}$.

  \textbf{Step 5: Terminal measurement.}  Measure the output qubit.  This
  constitutes the single superposition decay event.  By
  Theorem~\ref{thm:bqp}, the resulting problem is in $\bqp$.

  The total circuit size is $O(n \cdot \sqrt{k}) \leq O(n \cdot 2^{cn/2})$,
  which is polynomial for fixed $c < \tfrac{2 \log n}{n}$---i.e., for linear
  $f(n) = cn$ with small $c$.  This covers all algorithms with sub-exponential
  branch counts, which includes all deterministic and polynomially-branching
  non-deterministic linear-time algorithms. \qed
\end{proof}

\subsection{Linear Systems as Parallel DAGs}

\begin{theorem}[Linear Systems as Probabilistic DAGs]
  \label{thm:dag}
  Any system of $m$ linear equations in $n$ variables over $\Real$ (or
  $\Complex$) with a unique solution can be expressed as a probabilistic DAG
  $G = (V, E, p)$ whose parallel evaluation under \SD{} yields the solution
  in time $O(\log n)$ on a quantum system.
\end{theorem}

\begin{proof}[Proof Sketch]
  The system $A\mathbf{x} = \mathbf{b}$ (with $A \in \Real^{m \times n}$)
  has a unique solution $\mathbf{x}^* = A^{-1}\mathbf{b}$ (assuming $A$ is
  invertible).  The HHL algorithm~\cite{harrow2009} solves this system in
  $O(\log n \cdot \kappa^2)$ time on a quantum computer, where $\kappa$ is
  the condition number of $A$.  We model each variable as a DAG node and
  each dependency as a directed edge.  Gaussian elimination induces a
  topological order on the DAG; each elimination step is a parallel unitary
  operation on the corresponding qubit register.  The terminal measurement
  on each variable register constitutes a separate \SD{} event, and the
  results compose to give $\mathbf{x}^*$. \qed
\end{proof}

% ─────────────────────────────────────────────────────────────────────────────
\section{Rift: A Quantum-Classical Grammar Matcher}
\label{sec:rift}
% ─────────────────────────────────────────────────────────────────────────────

\subsection{Overview}

\textbf{Rift} (formally: \emph{RIFT --- a Flexible Translator}) is the
compiler front-end component of the OBINexus toolchain.  Its theoretical
grounding in the \SD{} framework arises from the observation that grammar
matching in a parser shares the structural properties of quantum superposition:
multiple production-rule paths may be simultaneously \emph{active} (pending)
until sufficient lookahead disambiguates them.  The decay event corresponds to
the moment the parser commits to a single parse path.

\subsection{Wave Postulates as Triangular Matching Postulates}

The wave postulate (Axiom~3) asserts that a quantum system in superposition
$\ket{\Psi} = c_1 \ket{\psi_1} + c_2 \ket{\psi_2}$ yields one of two outcomes
upon measurement.  We draw a direct structural analogy to a context-free
grammar $G = (V, \Sigma, R, S)$ with two competing production rules
$A \to \alpha \mid \beta$.  In both cases:

\begin{itemize}
  \item The system is in a \emph{pending} superposition of both alternatives
        until a discriminating token (the ``measurement'') is encountered.
  \item The collapse to a single rule (the ``decay'') is the parser's commit
        to one production.
  \item Destructive interference corresponds to the parser eliminating
        impossible parse paths via lookahead or semantic constraints.
\end{itemize}

\begin{definition}[Grammar Superposition State]
  Let $G = (V, \Sigma, R, S)$ be a context-free grammar and let $\pi$ be a
  partial parse of input $w = w_1 w_2 \cdots w_n$.  The \emph{grammar
  superposition state} at position $i$ is the set
  \[
    \Gamma_i \;=\; \bigl\{\, (A \to \alpha \bullet \beta,\, j)
                   \;\mid\; A \to \alpha\beta \in R,\;
                   w_j \cdots w_i \Rightarrow^* \alpha \,\bigr\}
  \]
  of all active Earley items.  The state $\Gamma_i$ is the grammar analogue
  of the superposition state vector: it encodes all simultaneously possible
  parse paths.
\end{definition}

\begin{definition}[Grammar Decay Event]
  A \emph{grammar decay event} at position $i$ occurs when
  $|\Gamma_i| = 1$, i.e., all but one parse path has been eliminated by
  lookahead.  This corresponds to the collapse of the grammar superposition
  state to a single production, exactly as measurement collapses
  $\ket{\Psi}$ to a single basis state.
\end{definition}

\subsection{Parallel Grammar Evaluation}

\begin{theorem}[Rift Parallel Parsing]
  \label{thm:rift}
  Let $G$ be an unambiguous context-free grammar and let $w$ be an input
  string of length $n$.  Rift evaluates all production rules of $G$ in
  \emph{parallel} by maintaining the full grammar superposition state
  $\Gamma_i$ at each position, collapsing (decaying) only when the
  lookahead forces $|\Gamma_i| = 1$.  The resulting parse is produced in
  $O(n^3)$ time in the worst case (Earley) and $O(n)$ for unambiguous
  grammars in Chomsky Normal Form, corresponding to a single-decay-event
  parse.
\end{theorem}

\begin{proof}[Proof Sketch]
  The Earley algorithm maintains all active items in $\Gamma_i$ at each
  position.  For an unambiguous grammar, the number of active items is
  bounded by a constant at each position after the initial scan, so the
  grammar superposition state decays to a singleton in $O(n)$ steps.  This
  is structurally identical to a $\bqp$ computation: the circuit evaluates
  all branches in parallel and the terminal measurement (grammar decay)
  yields the unique correct parse with probability 1. \qed
\end{proof}

\subsection{Stabilisation Principle}

The \SD{} framework introduces a \emph{stabilisation principle} as the
binding mechanism between the quantum measurement window and its cognitive
or computational substrate.

\begin{definition}[Stabilisation Principle]
  A system $S$ obeys the \emph{stabilisation principle} if, within the \SD{}
  window, all intermediate states are reversible (unitary), and the terminal
  measurement commits the system to a stable classical output that cannot be
  further modified within the same decay window.  Formally:
  \[
    \forall\, t < t_{\text{decay}}:\; U(t)\text{ is unitary and reversible.}
    \qquad
    \text{At } t_{\text{decay}}:\; \text{output is classical and irreversible.}
  \]
\end{definition}

In the Rift context, this means: all grammar transitions before the decay
event are reversible (backtracking is permitted); after the grammar decay
event, the committed parse tree is final.

% ─────────────────────────────────────────────────────────────────────────────
\section{Quantum Burst Protocol (QBP)}
\label{sec:qbp}
% ─────────────────────────────────────────────────────────────────────────────

\subsection{Motivation}

The \emph{Quantum Burst Protocol} (QBP) extends the \SD{} framework from
computation to communication.  Just as a quantum circuit collapses a
superposition of computation branches to a single classical output, a QBP
connection collapses a superposition of communication paths between two
endpoints to a single established channel.

\subsection{Formal Definition}

\begin{definition}[QBP Connection]
  A \emph{Quantum Burst Protocol connection} is a tuple
  \[
    \mathrm{QBP}(A, B) \;=\; \bigl(G,\; \SD(a,b),\; \Pi \bigr),
  \]
  where:
  \begin{itemize}
    \item $G = (E, N)$ is the quantum connection graph (Definition~3.1),
          with $A, B \in N$ the sender and receiver endpoints.
    \item $\SD(a,b)$ is the superposition decay window governing the
          connection establishment handshake.
    \item $\Pi = \{\pi_1, \pi_2, \ldots, \pi_k\}$ is the set of candidate
          communication paths (edges in $E$) held in superposition during
          the handshake phase.
  \end{itemize}
  The connection is \emph{established} (resolved) when the \SD{} window
  closes and a single path $\pi^* \in \Pi$ is selected; it is
  \emph{refused} (rejected) if no path survives the decay; and it is
  \emph{negotiating} (pending) during the handshake phase.
\end{definition}

\subsection{QBP Handshake as Superposition Decay}

\begin{proposition}[QBP Three-Phase Handshake]
  A QBP connection establishment proceeds in three phases corresponding
  exactly to the three states of the \SD{} window:
  \begin{enumerate}
    \item \textbf{Negotiating (pending):}  Both $A$ and $B$ maintain a
          superposition of all feasible communication paths $\Pi$.  No
          classical commitment is made.
    \item \textbf{Resolved:}  The \SD{} window closes; the path $\pi^*$
          with maximal probability amplitude is selected.  $A$ and $B$
          are classically connected via $\pi^*$.
    \item \textbf{Rejected:}  The \SD{} window closes with all path
          amplitudes below the threshold $b(n)$; no connection is established.
  \end{enumerate}
\end{proposition}

\begin{remark}
  The QBP handshake is a quantum analogue of the classical TCP three-way
  handshake (SYN / SYN-ACK / ACK), with the key difference that in QBP,
  \emph{all candidate paths are simultaneously active} during the negotiating
  phase, and the path selection is determined by a single quantum measurement
  (the decay event) rather than sequential message exchange.
\end{remark}

\subsection{Spectral Graph Representation}

The candidate path set $\Pi$ is naturally represented as a \emph{spectral
graph}:

\begin{definition}[Spectral Connection Graph]
  The \emph{spectral connection graph} of a QBP connection is the directed
  graph $\hat{G} = (V, E, \lambda)$ where $\lambda : E \to [0,1]$ assigns
  each edge a probability amplitude (eigenvalue weight) such that
  $\sum_{e \in \Pi} \lambda(e)^2 = 1$.  The selected path $\pi^*$ is the
  eigenvector corresponding to the dominant eigenvalue after the decay
  measurement.
\end{definition}

\begin{example}[Internet Connection as QBP]
  Consider connecting a client $A$ to a server $B$ over a network with $k$
  possible routing paths.  In classical networking, these paths are
  evaluated sequentially (or via parallel probes with separate classical
  outcomes).  Under QBP, the network interface holds all $k$ paths in
  superposition:
  \[
    \ket{\Pi} = \frac{1}{\sqrt{k}} \sum_{i=1}^k \ket{\pi_i}.
  \]
  Quality-of-service metrics (latency, bandwidth, packet loss) act as the
  observable $\hat{O}$.  Measurement of $\hat{O}$ collapses $\ket{\Pi}$ to
  the optimal path $\ket{\pi^*}$ in a single decay event, achieving optimal
  routing in $O(1)$ measurement time (assuming the superposition can be
  prepared in $O(\log k)$ time).
\end{example}

% ─────────────────────────────────────────────────────────────────────────────
\section{MMUKO OS and the OBINexus Toolchain}
\label{sec:mmuko}
% ─────────────────────────────────────────────────────────────────────────────

\subsection{MMUKO OS}

\textbf{MMUKO OS} (M for Mike, Q for Bella [Queen], K for Kilo, O for Oscar)
is the operating system substrate for the OBINexus quantum-classical computing
framework.  The name encodes its philosophical foundation: \emph{``the spirit
of good and evil that connects nothing and everything, that binds all of us and
all of them.''} This duality---nothing and everything, good and evil---maps
directly onto the binary superposition state $\ket{0}$ and $\ket{1}$ and their
linear combinations.

MMUKO OS is a cloud-based Windows system that serves as the classical
substrate on which quantum algorithms are emulated prior to the availability
of dedicated quantum hardware.  Its design is guided by three principles:

\begin{enumerate}
  \item \textbf{Read-Write symmetry:}  All memory operations respect the RWX
        model (Section~\ref{sec:rwx}), ensuring that write-heavy workloads
        collapse atomically to coherent read outcomes.
  \item \textbf{Superposition-aware scheduling:}  The OS scheduler treats
        parallel threads as superposed branches of a single computation,
        deferring classical commitment until the \SD{} window closes.
  \item \textbf{Decay-triggered synchronisation:}  Inter-process
        communication is governed by QBP handshakes
        (Section~\ref{sec:qbp}), replacing traditional lock-based
        synchronisation with measurement-based collapse events.
\end{enumerate}

\subsection{Toolchain Stack}

The OBINexus toolchain implements the \SD{} framework across the full
compilation pipeline:

\begin{center}
\renewcommand{\arraystretch}{1.4}
\begin{tabular}{@{}>{\ttfamily}p{3cm}p{9cm}@{}}
  \toprule
  \normalfont\textbf{Component} & \textbf{Role in \SD{} Framework} \\
  \midrule
  riftlang.exe   & Source language compiler; implements grammar superposition
                   (Section~\ref{sec:rift}) to parse Rift source files in
                   parallel. \\
  .so.a          & Shared object / static archive; intermediate representation
                   holding the superposed (pre-decay) compilation state. \\
  rift.exe       & Linked executable; represents the post-decay committed
                   classical binary output. \\
  gosilang       & Target runtime language binding; maps \SD{} primitives to
                   hardware-level instructions. \\
  nlink          & Build linker / orchestrator; manages the RWX transaction
                   across compilation units. \\
  polybuild      & Polyglot build system; coordinates multiple \SD{} decay
                   events across heterogeneous language targets. \\
  \bottomrule
\end{tabular}
\end{center}

The pipeline can be summarised as:
\[
  \texttt{riftlang.exe} \;\to\; \texttt{.so.a} \;\to\; \texttt{rift.exe}
  \;\to\; \texttt{gosilang} \;\xrightarrow{\;\texttt{nlink / polybuild}\;}\;
  \text{deployed binary}.
\]

% ─────────────────────────────────────────────────────────────────────────────
\section{Conclusion and Future Work}
\label{sec:conclusion}
% ─────────────────────────────────────────────────────────────────────────────

We have formalised the concept of \emph{Superposition Decay} as a
measurement window bridging quantum and classical computation.  The main
results are:

\begin{itemize}
  \item The \SD{} window is a three-valued construct (resolved / rejected /
        pending) that generalises the binary output of a classical TM and
        aligns precisely with the promise-problem formulation of quantum
        complexity classes.
  \item The RWX memory model provides an operational interpretation of the
        decay window: $n$ write operations and $k$ read operations within
        the window collapse atomically to a single coherent read.
  \item $\bqp$ is exactly the class of promise problems solvable within a
        single \SD{} event with a probability gap of $\tfrac{1}{3}$.
  \item Classical algorithms bounded by linear time complexity---deterministic
        or non-deterministic---can be simulated in parallel within a single
        \SD{} window by a polynomial-size quantum circuit.
  \item Systems of linear equations can be expressed as probabilistic DAGs
        whose parallel evaluation under \SD{} yields solutions in
        $O(\log n)$ quantum time.
  \item Rift's parser exploits grammar superposition to evaluate all
        production rules in parallel, collapsing to a unique parse tree at a
        grammar decay event governed by the stabilisation principle.
  \item The Quantum Burst Protocol (QBP) extends \SD{} to networking: all
        candidate communication paths are held in superposition during
        handshake, with optimal path selection achieved by a single spectral
        decay measurement.
  \item The MMUKO OS / OBINexus toolchain implements all of the above across
        the full pipeline: \texttt{riftlang.exe} $\to$ \texttt{.so.a} $\to$
        \texttt{rift.exe} $\to$ \texttt{gosilang}, orchestrated by
        \texttt{nlink} and \texttt{polybuild}.
\end{itemize}

\subsection{Future Work}

The following directions are planned within the OBINexus research programme:

\begin{enumerate}
  \item \textbf{Hardware realisation.}  Implementing the \SD{} framework on
        near-term quantum hardware (superconducting qubits, trapped ions) via
        the MMUKO OS / \texttt{riftlang} compiler toolchain.

  \item \textbf{Multi-decay protocols for \pp{} and \pspace.}  Designing
        iterated \SD{} protocols that amplify probability across multiple
        decay events to address problems beyond $\bqp$.

  \item \textbf{QBP networking stack.}  Implementing QBP as a full
        networking protocol atop the MMUKO OS kernel, replacing TCP
        handshakes with spectral decay measurements for latency-optimal
        routing.

  \item \textbf{Rift compiler validation.}  Formal verification that the
        Rift grammar decay event produces parse trees equivalent to those of
        a classical Earley parser, with empirical benchmarks on
        \texttt{gosilang} target programs.

  \item \textbf{Cognitive binding.}  Exploring the stabilisation principle
        as a model for attention and commitment in cognitive architectures:
        the \SD{} window as a formal model for the moment of decision in
        neural-symbolic systems.
\end{enumerate}

% ─────────────────────────────────────────────────────────────────────────────
\begin{thebibliography}{9}

\bibitem{grover1996}
  L.~K. Grover,
  ``A fast quantum mechanical algorithm for database search,''
  \emph{Proceedings of the 28th Annual ACM Symposium on Theory of Computing
  (STOC)}, pp.\ 212--219, 1996.

\bibitem{harrow2009}
  A.~W. Harrow, A.~Hassidim, and S.~Lloyd,
  ``Quantum algorithm for linear systems of equations,''
  \emph{Physical Review Letters}, vol.\ 103, no.\ 15, p.\ 150502, 2009.

\bibitem{nielsen2010}
  M.~A. Nielsen and I.~L. Chuang,
  \emph{Quantum Computation and Quantum Information},
  10th anniversary ed., Cambridge University Press, 2010.

\bibitem{watrous2009}
  J.~Watrous,
  ``Quantum computational complexity,''
  in \emph{Encyclopedia of Complexity and Systems Science},
  Springer, 2009.

\end{thebibliography}

\end{document}
